% !TEX root = ../main.tex

\chapter{Resultados}
\label{chap:resultados}

La metodología y los corpus de validación presentados en el capítulo anterior establecen las condiciones bajo las cuales se evalúa experimentalmente la solución implementada. En este capítulo se muestran los resultados obtenidos al aplicar este procedimiento sobre los programas de cada corpus, siendo \texttt{Maybe} y \texttt{Dist.T Rational} las mónadas sobre las que se valida la solución. Primeramente se presentarán los resultados de forma agrupada, de modo que sea posible observar el comportamiento general de todas las variantes sometidas a validación.

Luego de mostrar los resultados generales, se analizan los casos destacables del corpus probabilístico. Para cada uno de ellos se presenta el código fuente que lo conforma junto con su grafo de precedencia RAW asociado, lo que permite relacionar las dependencias presentes en el código con la estructura construida por la extensión. Posteriormente se examinan los programas de la familia \texttt{cost-selection} para estudiar el comportamiento de la selección del mejor plan de ejecución ante distintos perfiles de costo. Finalmente, se describen los resultados del control negativo con \texttt{IO} y se discute el alcance de la evidencia obtenida.

% !TEX root = ../../main.tex

\section{Resultados generales}
\label{sec:resultados-generales}

La metodología de validación de la solución mostrado en la Figura~\ref{fig:validacion-flujo} se aplicó a 93 programas que no correspondían a controles de activación: 44 del corpus de \texttt{Maybe} y 49 del corpus probabilístico. En todos ellos, la variante de \texttt{ApplicativeDo}, la selección automática y cada selección forzada conservaron tanto el código de salida como la salida observable del programa original. En particular, los casos de \texttt{Dist.T Rational} produjeron exactamente la misma distribución después de aplicar \texttt{Dist.norm}, incluidos aquellos con pesos no uniformes, resultados repetidos, soportes dependientes y condicionamiento. Por tanto, todos los casos evaluados respondieron favorablemente \textbf{PV1} dentro del criterio observacional establecido.

En los mismos 93 programas, el costo de cada candidato compilado individualmente permitió recalcular el mínimo, su primer índice y la cantidad de empates. En todos los casos, estos valores coincidieron con las métricas emitidas por el Renamer y el candidato elegido automáticamente correspondió al primer plan de costo mínimo entre los reordenamientos válidos generados. Este resultado responde favorablemente \textbf{PV2}. Además, como el orden fuente siempre pertenece al conjunto de candidatos, la selección automática nunca produjo un costo mayor que el obtenido por \texttt{ApplicativeDo} sobre ese mismo orden.

Los seis programas restantes asociados al control de activación de la extensión respondieron \textbf{PV3} y también produjeron el comportamiento esperado. En ambos corpus, los cuatro programas con bloques \texttt{do} ordinarios mantuvieron el marcador no activado, mostrado en el Renamer como \texttt{commutative-do = False}, incluso cuando \texttt{QualifiedDo} estaba habilitado, la interfaz conmutativa estaba importada y se había declarado la instancia correspondiente. Los dos bloques \texttt{do} escritos mediante \texttt{CD.do} produjeron \texttt{commutative-do = True} a travez de las trazas del Renamer, cada uno generando seis candidatos, incluyendo el orden original. En consecuencia, el uso completo de la solución requiere \texttt{ApplicativeDo}, \texttt{QualifiedDo}, la interfaz calificada, una instancia \texttt{CommutativeMonad} y el marcador local \texttt{CD.do} en el bloque que quiera ser declarado como conmutativo. La instancia expresa el contrato de conmutatividad y se comprueba durante el tipado, mientras que el marcador resoluble es la condición que distingue directamente ambas rutas dentro del Renamer.

La cantidad de casos que reducen estrictamente el costo de \texttt{ApplicativeDo} no se utiliza como indicador agregado de efectividad. La mayor parte del corpus fue construida para comprobar preservación observable, formación de dependencias, recolección de binders o selección correcta, y no para disponer deliberadamente un orden fuente desfavorable. Por ello, es esperable que una cadena de statements admita un único orden, que varios reordenamientos válidos tengan el mismo costo o que \texttt{ApplicativeDo} ya alcance el costo mínimo. Los cinco programas de cada corpus que obtuvieron mejores resultados respecto a \texttt{ApplicativeDo} fueron tres casos de la familia \texttt{cost-selection}, el caso asociado a dos cadenas de statements independientes y el caso en el que se usa \texttt{let} independiente.

Establecidos estos resultados globales, las secciones siguientes presentan casos representativos de las familias evaluadas. El objetivo es mostrar cómo los fenómenos presentes en el código fuente se reflejan en el grafo de precedencia RAW y cómo esa estructura determina los reordenamientos y planes que puede considerar la extensión.

% !TEX root = ../../main.tex

\section{Casos destacables}
\label{sec:resultados-casos-destacables}

Ya declarados los resultados globales de los dos corpus evaluados, en esta sección se examinarán casos representativos de las familias de casos descritas en el capítulo anterior, las cuales se encuentran en las Tablas~\ref{tab:validacion-familias-compartidas} y \ref{tab:validacion-familias-especificas}. Se utilizan las variantes de \texttt{Dist.T Rational} para mantener el análisis dentro del dominio probabilístico principal de la memoria; las variantes estructurales equivalentes de \texttt{Maybe} produjeron los mismos grafos y métricas.

\subsection{Formas de dependencia}

La familia \texttt{dependency-shapes} permite contrastar la cantidad de reordenamientos válidos generados por la solución con cantidades conocidas que se obtienen en función de la forma del grafo de precedencia RAW. Los cuatro casos siguientes abarcan desde una secuencia de statements no reordenable hasta un programa con un espacio de exploración de 70 reordenamientos válidos.

\subsubsection{Cadena completa}

El caso \texttt{chain-no-reorder} consiste en un bloque \texttt{do} en el que cada distribución utiliza el binder definido por el statement anterior. Esta forma produce una cadena completa de dependencias y fija el orden de todos los statements. A continuación, en el Código~\ref{lst:resultados-cadena} y la Figura~\ref{fig:resultados-grafo-cadena} se muestran el programa probabilístico asociado a este caso y el grafo de precedencia RAW construido por la extensión, respectivamente.

\begin{lstlisting}[style=haskell,caption={Programa probabilístico asociado al caso \texttt{chain-no-reorder}.},label={lst:resultados-cadena}]
CD.do
  x1 <- Dist.uniform [1, 2]
  x2 <- Dist.uniform [x1 + 10, x1 + 20]
  x3 <- Dist.uniform [x2 + 100, x2 + 200]
  x4 <- Dist.uniform [x3 + 1000, x3 + 2000]
  CD.return (x1, x2, x3, x4)
\end{lstlisting}

\begin{figure}[ht]
\centering
\begin{PrecedenceGraph}
    \PrecedenceNode{s1}{x1}{$s_1$}
    \PrecedenceNode[right=of s1]{s2}{x2}{$s_2$}
    \PrecedenceNode[right=of s2]{s3}{x3}{$s_3$}
    \PrecedenceNode[right=of s3]{s4}{x4}{$s_4$}
    \PrecedenceEdge{s1}{s2}
    \PrecedenceEdge{s2}{s3}
    \PrecedenceEdge{s3}{s4}
\end{PrecedenceGraph}
\caption[Grafo RAW del caso \texttt{chain-no-reorder}.]{Grafo de precedencia RAW construido por la solución para el caso \texttt{chain-no-reorder}.}
\label{fig:resultados-grafo-cadena}
\end{figure}

Como muestra la Figura~\ref{fig:resultados-grafo-cadena}, las aristas $s_1\to s_2\to s_3\to s_4$ fijan por completo el orden de los cuatro statements. Por esta razón, la solución solo generó el candidato correspondiente al orden fuente y no encontró órdenes alternativos. La ejecución secuencial, \texttt{ApplicativeDo} y la solución obtuvieron un costo de 4. Este resultado confirma que las dependencias RAW restringen la generación de reordenamientos cuando todos los statements forman una cadena.

\subsubsection{Diamante de seis statements}

El caso \texttt{diamond} consiste en un bloque \texttt{do} que comienza con una distribución común, continúa con dos cadenas de distribuciones independientes entre sí y finaliza con una distribución que utiliza los resultados de ambas. Esta forma produce dos ramas que se separan a partir de \texttt{root} y convergen en \texttt{join}. A continuación, en el Código~\ref{lst:resultados-diamante} y la Figura~\ref{fig:resultados-grafo-diamante} se muestran el programa probabilístico asociado a este caso y el grafo de precedencia RAW construido por la extensión, respectivamente.

\begin{lstlisting}[style=haskell,caption={Programa probabilístico asociado al caso \texttt{diamond}.},label={lst:resultados-diamante}]
CD.do
  root <- Dist.uniform [1, 2]
  a1 <- Dist.uniform [root + 10, root + 20]
  b1 <- Dist.uniform [root + 100, root + 200]
  a2 <- Dist.uniform [a1 + 1000, a1 + 2000]
  b2 <- Dist.uniform [b1 + 10000, b1 + 20000]
  join <- Dist.uniform [a2 + b2 + 100000, a2 + b2 + 200000]
  CD.return (root, a1, a2, b1, b2, join)
\end{lstlisting}

\begin{figure}[ht]
\centering
\begingroup
\setgraphnodesize{1.55cm}{1.0cm}{8pt}
\begin{PrecedenceGraph}[node distance=0.45cm and 0.35cm]
    \PrecedenceNode{s1}{root}{$s_1$}
    \PrecedenceNode[right=of s1]{s2}{a1}{$s_2$}
    \PrecedenceNode[right=of s2]{s3}{b1}{$s_3$}
    \PrecedenceNode[right=of s3]{s4}{a2}{$s_4$}
    \PrecedenceNode[right=of s4]{s5}{b2}{$s_5$}
    \PrecedenceNode[right=of s5]{s6}{join}{$s_6$}
    \PrecedenceEdge{s1}{s2}
    \PrecedenceEdge[bend right=48]{s1}{s3}
    \PrecedenceEdge[bend left=35]{s2}{s4}
    \PrecedenceEdge[bend right=45]{s3}{s5}
    \PrecedenceEdge[bend left=40]{s4}{s6}
    \PrecedenceEdge{s5}{s6}
\end{PrecedenceGraph}
\endgroup
\caption[Grafo RAW del caso \texttt{diamond}.]{Grafo de precedencia RAW construido por la solución para el caso \texttt{diamond}.}
\label{fig:resultados-grafo-diamante}
\end{figure}

La Figura~\ref{fig:resultados-grafo-diamante} muestra un grafo que contiene las cadenas de aristas $s_1\to s_2\to s_4\to s_6$ y $s_1\to s_3\to s_5\to s_6$. Los statements de una rama pueden intercalarse con los de la otra siempre que cada cadena conserve su orden, por lo que el grafo admite seis órdenes válidos, cuatro de ellos con costo mínimo. La ejecución secuencial obtuvo un costo de 6, mientras que \texttt{ApplicativeDo} y la solución obtuvieron un costo de 4. Como \texttt{ApplicativeDo} ya construyó un plan mínimo sobre el orden fuente, la solución conservó ese primer candidato.

\subsubsection{Grafo sin aristas}

El caso \texttt{no-deps} consiste en un bloque \texttt{do} con cuatro distribuciones que no leen binders locales producidos por otros statements. Esta independencia produce un grafo de precedencia RAW sin aristas. A continuación, en el Código~\ref{lst:resultados-grafo-vacio} y la Figura~\ref{fig:resultados-grafo-vacio} se muestran el programa probabilístico asociado a este caso y el grafo de precedencia RAW construido por la extensión, respectivamente.

\begin{lstlisting}[style=haskell,caption={Programa probabilístico asociado al caso \texttt{no-deps}.},label={lst:resultados-grafo-vacio}]
CD.do
  x1 <- Dist.uniform [1, 2]
  x2 <- Dist.uniform [10, 20]
  x3 <- Dist.uniform [100, 200]
  x4 <- Dist.uniform [1000, 2000]
  CD.return (x1, x2, x3, x4)
\end{lstlisting}

\begin{figure}[ht]
\centering
\begin{PrecedenceGraph}
    \PrecedenceNode{s1}{x1}{$s_1$}
    \PrecedenceNode[right=of s1]{s2}{x2}{$s_2$}
    \PrecedenceNode[right=of s2]{s3}{x3}{$s_3$}
    \PrecedenceNode[right=of s3]{s4}{x4}{$s_4$}
\end{PrecedenceGraph}
\caption[Grafo RAW del caso \texttt{no-deps}.]{Grafo de precedencia RAW construido por la solución para el caso \texttt{no-deps}.}
\label{fig:resultados-grafo-vacio}
\end{figure}

El grafo de la Figura~\ref{fig:resultados-grafo-vacio} no contiene aristas, de modo que los cuatro statements pueden aparecer en cualquier orden. El grafo admite 24 órdenes válidos y todos alcanzaron el costo mínimo. La ejecución secuencial obtuvo un costo de 4, mientras que \texttt{ApplicativeDo} y la solución obtuvieron un costo de 1 al agrupar las cuatro distribuciones independientes. Como todos los candidatos tuvieron el mismo costo, la solución conservó el orden fuente.

\subsubsection{Dos cadenas independientes}

El caso \texttt{two-chains} consiste en un bloque \texttt{do} formado por dos cadenas independientes de cuatro distribuciones cada una. En el orden fuente, el último statement de cada cadena aparece separado de los tres anteriores, lo que dificulta que \texttt{ApplicativeDo} agrupe ambas cadenas completas. A continuación, en el Código~\ref{lst:resultados-dos-cadenas} y la Figura~\ref{fig:resultados-grafo-dos-cadenas} se muestran el programa probabilístico asociado a este caso y el grafo de precedencia RAW construido por la extensión, respectivamente.

\begin{lstlisting}[style=haskell,caption={Programa probabilístico asociado al caso \texttt{two-chains}.},label={lst:resultados-dos-cadenas}]
CD.do
  a1 <- Dist.uniform [1, 2]
  a2 <- Dist.uniform [a1 + 100, a1 + 200]
  a3 <- Dist.uniform [a2 + 10000, a2 + 20000]
  b1 <- Dist.uniform [10, 20]
  b2 <- Dist.uniform [b1 + 1000, b1 + 2000]
  b3 <- Dist.uniform [b2 + 100000, b2 + 200000]
  a4 <- Dist.uniform [a3 + 1000000, a3 + 2000000]
  b4 <- Dist.uniform [b3 + 10000000, b3 + 20000000]
  CD.return (a1, a2, a3, a4, b1, b2, b3, b4)
\end{lstlisting}

\begin{figure}[ht]
\centering
\begingroup
\setgraphnodesize{1.55cm}{0.85cm}{8pt}
\begin{PrecedenceGraph}[node distance=0.35cm and 0.2cm]
    \PrecedenceNode{s1}{a1}{$s_1$}
    \PrecedenceNode[right=of s1]{s2}{a2}{$s_2$}
    \PrecedenceNode[right=of s2]{s3}{a3}{$s_3$}
    \PrecedenceNode[right=of s3]{s4}{b1}{$s_4$}
    \PrecedenceNode[right=of s4]{s5}{b2}{$s_5$}
    \PrecedenceNode[right=of s5]{s6}{b3}{$s_6$}
    \PrecedenceNode[right=of s6]{s7}{a4}{$s_7$}
    \PrecedenceNode[right=of s7]{s8}{b4}{$s_8$}
    \PrecedenceEdge{s1}{s2}
    \PrecedenceEdge{s2}{s3}
    \PrecedenceEdge[bend left=35]{s3}{s7}
    \PrecedenceEdge{s4}{s5}
    \PrecedenceEdge{s5}{s6}
    \PrecedenceEdge[bend right=60]{s6}{s8}
\end{PrecedenceGraph}
\endgroup
\caption[Grafo RAW del caso \texttt{two-chains}.]{Grafo de precedencia RAW construido por la solución para el caso \texttt{two-chains}.}
\label{fig:resultados-grafo-dos-cadenas}
\end{figure}

Como se muestra en la Figura~\ref{fig:resultados-grafo-dos-cadenas}, el grafo contiene la cadena de aristas $s_1\to s_2\to s_3\to s_7$ junto con $s_4\to s_5\to s_6\to s_8$. Como no existen dependencias entre ambas cadenas, sus statements pueden intercalarse de 70 formas, 30 de las cuales alcanzaron el costo mínimo. La ejecución secuencial obtuvo un costo de 8 y \texttt{ApplicativeDo} redujo ese valor a 6. La solución reunió los statements de cada cadena y obtuvo un costo de 4, lo que muestra que un reordenamiento puede exponer agrupaciones aplicativas que no eran visibles en el orden fuente.

\subsection{Dependencias introducidas por statements let}

El caso \texttt{bind-depends-on-let} consiste en un bloque \texttt{do} cuyo primer statement define \texttt{x1} mediante \texttt{let}, mientras que el segundo utiliza ese identificador para construir una distribución y el tercero permanece independiente. Esta forma permite observar una dependencia RAW cuyo productor no es un binder monádico. A continuación, en el Código~\ref{lst:resultados-bind-depends-on-let} y la Figura~\ref{fig:resultados-grafo-bind-depends-on-let} se muestran el programa probabilístico asociado a este caso y el grafo de precedencia RAW construido por la extensión, respectivamente.

\begin{lstlisting}[style=haskell,caption={Programa probabilístico asociado al caso \texttt{bind-depends-on-let}.},label={lst:resultados-bind-depends-on-let}]
CD.do
  let x1 = 10
  x2 <- Dist.uniform [x1 + 1, x1 + 2]
  x3 <- Dist.uniform [100, 200]
  CD.return (x1, x2, x3)
\end{lstlisting}

\begin{figure}[ht]
\centering
\begin{PrecedenceGraph}[node distance=0.5cm and 0.6cm]
    \PrecedenceNode{s1}{let x1}{$s_1$}
    \PrecedenceNode[right=of s1]{s2}{x2}{$s_2$}
    \PrecedenceNode[right=of s2]{s3}{x3}{$s_3$}
    \PrecedenceEdge{s1}{s2}
\end{PrecedenceGraph}
\caption[Grafo RAW del caso \texttt{bind-depends-on-let}.]{Grafo de precedencia RAW construido por la solución para el caso \texttt{bind-depends-on-let}.}
\label{fig:resultados-grafo-bind-depends-on-let}
\end{figure}

La Figura~\ref{fig:resultados-grafo-bind-depends-on-let} muestra un grafo de precedencia que contiene la arista $s_1\to s_2$, cuyo identificador usado es \texttt{x1}, mientras que $s_3$ no presenta dependencias con los otros statements. El grafo admite tres órdenes válidos y todos alcanzaron el costo mínimo. La ejecución secuencial obtuvo un costo de 3, mientras que \texttt{ApplicativeDo} y la solución obtuvieron un costo de 2 al agrupar el tercer statement con la secuencia formada por los dos primeros. Este resultado confirma que la extensión considera las declaraciones \texttt{let} como productoras de binders al construir las precedencias RAW.

\subsection{Binders definidos mediante patrones}

El caso \texttt{tuple-pattern} consiste en un bloque \texttt{do} cuyo primer statement utiliza un patrón de tupla para definir simultáneamente los binders \texttt{x1} y \texttt{x2}. Los dos statements siguientes consumen estos binders por separado, lo que permite observar cómo la extensión recolecta múltiples nombres desde un mismo patrón. A continuación, en el Código~\ref{lst:resultados-tuple-binder} y la Figura~\ref{fig:resultados-grafo-tuple-binder} se muestran el programa probabilístico asociado a este caso y el grafo de precedencia RAW construido por la extensión, respectivamente.

\begin{lstlisting}[style=haskell,caption={Programa probabilístico asociado al caso \texttt{tuple-pattern}.},label={lst:resultados-tuple-binder}]
CD.do
  (x1, x2) <- Dist.uniform [(1, 10), (2, 20)]
  x3 <- Dist.uniform [x1 + 1, x1 + 2]
  x4 <- Dist.uniform [x2 + 1, x2 + 2]
  CD.return (x1, x2, x3, x4)
\end{lstlisting}

\begin{figure}[ht]
\centering
\begingroup
\begin{PrecedenceGraph}[node distance=0.45cm and 0.55cm]
    \PrecedenceNode{s1}{(x1, x2)}{$s_1$}
    \PrecedenceNode[right=of s1]{s2}{x3}{$s_2$}
    \PrecedenceNode[right=of s2]{s3}{x4}{$s_3$}
    \PrecedenceEdge{s1}{s2}
    \PrecedenceEdge[bend left=35]{s1}{s3}
\end{PrecedenceGraph}
\endgroup
\caption[Grafo RAW del caso \texttt{tuple-pattern}.]{Grafo de precedencia RAW construido por la solución para el caso \texttt{tuple-pattern}.}
\label{fig:resultados-grafo-tuple-binder}
\end{figure}

El grafo de la Figura~\ref{fig:resultados-grafo-tuple-binder} contiene las aristas $s_1\to s_2$ y $s_1\to s_3$, cuyos identificadores consumidos son \texttt{x1} y \texttt{x2}, respectivamente. Como los dos statements posteriores no dependen entre sí, el grafo admite dos órdenes válidos y ambos alcanzaron el costo mínimo. La ejecución secuencial obtuvo un costo de 3, mientras que \texttt{ApplicativeDo} y la solución obtuvieron un costo de 2. Este resultado confirma que la extensión recolecta por separado los binders definidos dentro de un patrón de tupla.

\subsection{Distribuciones con parámetros dependientes}

El caso \texttt{choose-dependent} consiste en un bloque \texttt{do} en el que el valor de \texttt{x1} determina los pesos utilizados por \texttt{Dist.choose} para producir \texttt{x2}, mientras que \texttt{x3} permanece independiente. Esta forma permite observar una dependencia introducida dentro de una expresión condicional que modifica los parámetros de una distribución. A continuación, en el Código~\ref{lst:resultados-choose-dependent} y la Figura~\ref{fig:resultados-grafo-choose-dependent} se muestran el programa probabilístico asociado a este caso y el grafo de precedencia RAW construido por la extensión, respectivamente.

\begin{lstlisting}[style=haskell,caption={Programa probabilístico asociado al caso \texttt{choose-dependent}.},label={lst:resultados-choose-dependent}]
CD.do
  x1 <- Dist.uniform [1, 2]
  x2 <- if x1 == 1
        then Dist.choose (1 / 4) 10 20
        else Dist.choose (3 / 4) 10 20
  x3 <- Dist.uniform [100, 200]
  CD.return (x1, x2, x3)
\end{lstlisting}

\begin{figure}[ht]
\centering
\begin{PrecedenceGraph}[node distance=0.5cm and 0.6cm]
    \PrecedenceNode{s1}{x1}{$s_1$}
    \PrecedenceNode[right=of s1]{s2}{x2}{$s_2$}
    \PrecedenceNode[right=of s2]{s3}{x3}{$s_3$}
    \PrecedenceEdge{s1}{s2}
\end{PrecedenceGraph}
\caption[Grafo RAW del caso \texttt{choose-dependent}.]{Grafo de precedencia RAW construido por la solución para el caso \texttt{choose-dependent}.}
\label{fig:resultados-grafo-choose-dependent}
\end{figure}

La condición que selecciona los pesos lee \texttt{x1}, por lo que la Figura~\ref{fig:resultados-grafo-choose-dependent} contiene la arista $s_1\to s_2$, mientras que $s_3$ no presenta dependencias con los otros statements. El grafo admite tres órdenes válidos y todos alcanzaron el costo mínimo. La ejecución secuencial obtuvo un costo de 3, mientras que \texttt{ApplicativeDo} y la solución obtuvieron un costo de 2. Los tres órdenes conservaron la distribución normalizada y las probabilidades no uniformes definidas por cada rama, lo que muestra que la extensión detecta lecturas dentro de expresiones condicionales.

\subsection{Condicionamiento de una distribución dependiente}

El caso \texttt{filter-bind} consiste en un bloque \texttt{do} en el que el binder \texttt{x1} determina el soporte de la distribución entregada a \texttt{Dist.filter} en el segundo statement, mientras que el tercero permanece independiente. Esta forma representa el condicionamiento de una distribución que depende de un resultado producido previamente. A continuación, en el Código~\ref{lst:resultados-filter-bind} y la Figura~\ref{fig:resultados-grafo-filter-bind} se muestran el programa probabilístico asociado a este caso y el grafo de precedencia RAW construido por la extensión, respectivamente.

\begin{lstlisting}[style=haskell,caption={Programa probabilístico asociado al caso \texttt{filter-bind}.},label={lst:resultados-filter-bind}]
CD.do
  x1 <- Dist.uniform [1, 2]
  x2 <- Dist.filter even (Dist.uniform [x1, x1 + 1, x1 + 2])
  x3 <- Dist.uniform [100, 200]
  CD.return (x1, x2, x3)
\end{lstlisting}

\begin{figure}[ht]
\centering
\begin{PrecedenceGraph}[node distance=0.5cm and 0.6cm]
    \PrecedenceNode{s1}{x1}{$s_1$}
    \PrecedenceNode[right=of s1]{s2}{x2}{$s_2$}
    \PrecedenceNode[right=of s2]{s3}{x3}{$s_3$}
    \PrecedenceEdge{s1}{s2}
\end{PrecedenceGraph}
\caption[Grafo RAW del caso \texttt{filter-bind}.]{Grafo de precedencia RAW construido por la solución para el caso \texttt{filter-bind}.}
\label{fig:resultados-grafo-filter-bind}
\end{figure}

La lectura de \texttt{x1} dentro de la distribución filtrada produce la arista $s_1\to s_2$ de la Figura~\ref{fig:resultados-grafo-filter-bind}, mientras que $s_3$ no presenta dependencias con los otros statements. El grafo admite tres órdenes válidos y todos alcanzaron el costo mínimo. La ejecución secuencial obtuvo un costo de 3, mientras que \texttt{ApplicativeDo} y la solución obtuvieron un costo de 2. Los tres órdenes conservaron el soporte filtrado y sus probabilidades normalizadas, lo que confirma que la extensión mantiene la precedencia cuando una lectura aparece dentro de una operación de condicionamiento.

\subsection{Identidades únicas frente a shadowing}

El caso \texttt{read-after-rebind} consiste en un bloque \texttt{do} que reutiliza el identificador textual \texttt{x} en los dos primeros statements y luego lee ese identificador en el tercero para definir \texttt{y}. La segunda vinculación oculta a la primera, por lo que esta lectura debe depender del binder más reciente y no del \texttt{x} original. A continuación, en el Código~\ref{lst:resultados-shadowing} y la Figura~\ref{fig:resultados-grafo-shadowing} se muestran el programa probabilístico asociado a este caso y el grafo de precedencia RAW construido por la extensión, respectivamente.

\begin{lstlisting}[style=haskell,caption={Programa probabilístico asociado al caso \texttt{read-after-rebind}.},label={lst:resultados-shadowing}]
CD.do
  x <- Dist.uniform [1, 2]
  x <- Dist.uniform [10, 20]
  y <- Dist.uniform [x + 100, x + 200]
  CD.return (x, y)
\end{lstlisting}

\begin{figure}[ht]
\centering
\begingroup
\begin{PrecedenceGraph}[node distance=0.45cm and 0.5cm]
    \PrecedenceNode{s1}{x\_a1Y4}{$s_1$}
    \PrecedenceNode[right=of s1]{s2}{x\_a1Y5}{$s_2$}
    \PrecedenceNode[right=of s2]{s3}{y\_a1Y6}{$s_3$}
    \PrecedenceEdge{s2}{s3}
\end{PrecedenceGraph}
\endgroup
\caption[Grafo RAW del caso \texttt{read-after-rebind}.]{Grafo de precedencia RAW construido por la solución para el caso \texttt{read-after-rebind} utilizando los identificadores ya renombrados por el Renamer.}
\label{fig:resultados-grafo-shadowing}
\end{figure}

La Figura~\ref{fig:resultados-grafo-shadowing} contiene la arista $s_2\to s_3$, cuya lectura es \texttt{x\_a1Y5}. No existe una arista desde $s_1$, porque las identidades \texttt{x\_a1Y4} y \texttt{x\_a1Y5} representan binders distintos y la lectura del tercer statement fue asociada con el segundo. Por tanto, el primer statement permanece independiente de la cadena formada por los dos siguientes. El grafo admite tres órdenes válidos y todos alcanzaron el costo mínimo. La ejecución secuencial obtuvo un costo de 3, mientras que \texttt{ApplicativeDo} y la solución obtuvieron un costo de 2. Los tres órdenes conservaron la distribución normalizada, lo que confirma que la extensión respeta el shadowing y construye la dependencia RAW a partir de la identidad actualmente visible.

% !TEX root = ../../main.tex

\section{Selección por costo}
\label{sec:resultados-costos}

La familia de casos \texttt{cost-selection} fue diseñada para observar la política de selección del plan de ejecución de menor costo bajo cuatro perfiles distintos. La Tabla~\ref{tab:resultados-perfiles-costo} separa las tres métricas de costo y los dos conteos relevantes. \textit{Sec} corresponde al costo base de ejecutar todos los statements en secuencia; \textit{ADo}, al plan obtenido por \texttt{ApplicativeDo} sobre el orden fuente; y \textit{Extensión}, al menor costo encontrado después de evaluar los reordenamientos válidos.

\begin{table}[ht]
\centering
\begingroup
\setlength{\TableColumnPadding}{5pt}
\TableSetup
\TableFrame{%
\begin{tabular}{p{4.1cm}ccccc}
\rowcolor{tableHeaderBg}
\TableHead{Caso} & \TableHead{Sec} & \TableHead{ADo} & \TableHead{Extensión} & \TableHead{\shortstack{Reordenamientos\\totales}} & \TableHead{\shortstack{Reordenamientos\\costo mínimo}} \\
\texttt{all-same-cost} & 4 & 2 & 2 & 6 & 6 \\
\texttt{original-not-minimal} & 4 & 3 & 2 & 5 & 4 \\
\texttt{unique-minimum} & 5 & 5 & 4 & 5 & 1 \\
\texttt{many-minimums} & 6 & 4 & 3 & 20 & 14 \\
\end{tabular}%
}
\endgroup
\caption{Costos y cantidades de reordenamientos de la familia \texttt{cost-selection}.}
\label{tab:resultados-perfiles-costo}
\end{table}

\subsection{Todos los candidatos con el mismo costo}

El caso \texttt{all-same-cost} consiste en un bloque \texttt{do} cuyo primer statement define \texttt{x1} y cuyos tres statements siguientes utilizan ese binder sin depender entre sí. Esta forma produce una raíz común con tres sucesores independientes y permite observar qué ocurre cuando todos los reordenamientos válidos tienen el mismo costo. A continuación, en el Código~\ref{lst:resultados-all-same-cost} y la Figura~\ref{fig:resultados-grafo-all-same-cost} se muestran el programa probabilístico asociado a este caso y el grafo de precedencia RAW construido por la extensión, respectivamente.

\begin{lstlisting}[style=haskell,caption={Programa probabilístico asociado al caso \texttt{all-same-cost}.},label={lst:resultados-all-same-cost}]
CD.do
  x1 <- Dist.uniform [1, 2]
  x2 <- Dist.uniform [x1 + 10, x1 + 20]
  x3 <- Dist.uniform [x1 + 100, x1 + 200]
  x4 <- Dist.uniform [x1 + 1000, x1 + 2000]
  CD.return (x1, x2, x3, x4)
\end{lstlisting}

\begin{figure}[ht]
\centering
\begingroup
\setgraphnodesize{1.7cm}{1.0cm}{8pt}
\begin{PrecedenceGraph}[node distance=0.4cm and 0.45cm]
    \PrecedenceNode{s1}{x1}{$s_1$}
    \PrecedenceNode[right=of s1]{s2}{x2}{$s_2$}
    \PrecedenceNode[right=of s2]{s3}{x3}{$s_3$}
    \PrecedenceNode[right=of s3]{s4}{x4}{$s_4$}
    \PrecedenceEdge{s1}{s2}
    \PrecedenceEdge[bend left=28]{s1}{s3}
    \PrecedenceEdge[bend left=45]{s1}{s4}
\end{PrecedenceGraph}
\endgroup
\caption[Grafo RAW del caso \texttt{all-same-cost}.]{Grafo de precedencia RAW construido por la solución para el caso \texttt{all-same-cost}.}
\label{fig:resultados-grafo-all-same-cost}
\end{figure}

La Figura~\ref{fig:resultados-grafo-all-same-cost} contiene las aristas $s_1\to s_2$, $s_1\to s_3$ y $s_1\to s_4$. Una vez ejecutado el primer statement, los otros tres pueden aparecer en cualquier orden, por lo que el grafo admite seis órdenes válidos y todos alcanzaron el costo mínimo. La ejecución secuencial obtuvo un costo de 4, mientras que \texttt{ApplicativeDo} y la solución obtuvieron un costo de 2 al agrupar los tres sucesores independientes. Como todos los candidatos tuvieron el mismo costo, la solución conservó de manera estable el orden fuente.

\subsection{Orden original no mínimo}

El caso \texttt{original-not-minimal} consiste en un bloque \texttt{do} en el que \texttt{x2} depende de \texttt{x1}, \texttt{x3} permanece independiente y \texttt{x4} utiliza tanto \texttt{x1} como \texttt{x3}. En el orden fuente, el statement que define \texttt{x2} aparece antes que la distribución independiente, lo que impide que \texttt{ApplicativeDo} exponga directamente las dos etapas aplicativas del grafo. Esta forma permite observar si la solución abandona el orden original cuando existe un reordenamiento de menor costo. A continuación, en el Código~\ref{lst:resultados-original-not-minimal} y la Figura~\ref{fig:resultados-grafo-original-not-minimal} se muestran el programa probabilístico asociado a este caso y el grafo de precedencia RAW construido por la extensión, respectivamente.

\begin{lstlisting}[style=haskell,caption={Programa probabilístico asociado al caso \texttt{original-not-minimal}.},label={lst:resultados-original-not-minimal}]
CD.do
  x1 <- Dist.uniform [1, 2]
  x2 <- Dist.uniform [x1 + 10, x1 + 20]
  x3 <- Dist.uniform [100, 200]
  x4 <- Dist.uniform [x1 + x3, x1 + x3 + 1]
  CD.return (x1, x2, x3, x4)
\end{lstlisting}

\begin{figure}[ht]
\centering
\begingroup
\setgraphnodesize{1.7cm}{1.0cm}{8pt}
\begin{PrecedenceGraph}[node distance=0.4cm and 0.45cm]
    \PrecedenceNode{s1}{x1}{$s_1$}
    \PrecedenceNode[right=of s1]{s2}{x2}{$s_2$}
    \PrecedenceNode[right=of s2]{s3}{x3}{$s_3$}
    \PrecedenceNode[right=of s3]{s4}{x4}{$s_4$}
    \PrecedenceEdge{s1}{s2}
    \PrecedenceEdge[bend left=42]{s1}{s4}
    \PrecedenceEdge{s3}{s4}
\end{PrecedenceGraph}
\endgroup
\caption[Grafo RAW del caso \texttt{original-not-minimal}.]{Grafo de precedencia RAW construido por la solución para el caso \texttt{original-not-minimal}.}
\label{fig:resultados-grafo-original-not-minimal}
\end{figure}

La Figura~\ref{fig:resultados-grafo-original-not-minimal} contiene las aristas $s_1\to s_2$, $s_1\to s_4$ y $s_3\to s_4$. El grafo admite cinco órdenes válidos. La ejecución secuencial obtuvo un costo de 4 y \texttt{ApplicativeDo} obtuvo un costo de 3 sobre el orden fuente, que fue el único candidato con ese valor. Los otros cuatro órdenes alcanzaron un costo de 2; la solución seleccionó el primero de ellos, agrupando inicialmente los statements independientes $s_1$ y $s_3$ y luego los sucesores $s_2$ y $s_4$. Este resultado confirma que la política abandona el orden original cuando existe una alternativa de menor costo.

\subsection{Mínimo único}

El caso \texttt{unique-minimum} consiste en un bloque \texttt{do} que combina una cadena de cuatro statements con un statement independiente que utiliza un patrón de tupla. Aunque este último no presenta dependencias RAW, su patrón potencialmente estricto hace que \texttt{ApplicativeDo} conserve la secuencia con el statement siguiente cuando aparece antes del final del bloque. Esta forma produce un único reordenamiento de costo mínimo. A continuación, en el Código~\ref{lst:resultados-unique-minimum} y la Figura~\ref{fig:resultados-grafo-unique-minimum} se muestran el programa probabilístico asociado a este caso y el grafo de precedencia RAW construido por la extensión, respectivamente.

\begin{lstlisting}[style=haskell,caption={Programa probabilístico asociado al caso \texttt{unique-minimum}.},label={lst:resultados-unique-minimum}]
CD.do
  x1 <- Dist.uniform [1, 2]
  x2 <- Dist.uniform [x1 + 10, x1 + 20]
  x3 <- Dist.uniform [x2 + 100, x2 + 200]
  (x4, _) <- Dist.uniform [(40, 0), (50, 0)]
  x5 <- Dist.uniform [x3 + 1000, x3 + 2000]
  CD.return (x1, x2, x3, x4, x5)
\end{lstlisting}

\begin{figure}[ht]
\centering
\begingroup
\setgraphnodesize{1.5cm}{0.95cm}{8pt}
\begin{PrecedenceGraph}[node distance=0.4cm and 0.35cm]
    \PrecedenceNode{s1}{x1}{$s_1$}
    \PrecedenceNode[right=of s1]{s2}{x2}{$s_2$}
    \PrecedenceNode[right=of s2]{s3}{x3}{$s_3$}
    \PrecedenceNode[right=of s3]{s4}{(x4, \_)}{$s_4$}
    \PrecedenceNode[right=of s4]{s5}{x5}{$s_5$}
    \PrecedenceEdge{s1}{s2}
    \PrecedenceEdge{s2}{s3}
    \PrecedenceEdge[bend left=35]{s3}{s5}
\end{PrecedenceGraph}
\endgroup
\caption[Grafo RAW del caso \texttt{unique-minimum}.]{Grafo de precedencia RAW construido por la solución para el caso \texttt{unique-minimum}.}
\label{fig:resultados-grafo-unique-minimum}
\end{figure}

La Figura~\ref{fig:resultados-grafo-unique-minimum} contiene la cadena de aristas $s_1\to s_2\to s_3\to s_5$, mientras que $s_4$ permanece aislado porque su expresión no lee binders locales y \texttt{x4} no es utilizado por otro statement del grafo. Los cinco órdenes válidos corresponden a las posiciones en que puede insertarse $s_4$ respecto de la cadena. La ejecución secuencial y \texttt{ApplicativeDo} obtuvieron un costo de 5 sobre el orden fuente. Cuatro candidatos conservaron ese costo y solo el orden que ubica $s_4$ después de $s_5$ alcanzó un costo de 4, por lo que la solución seleccionó ese mínimo único y agrupó la cadena completa con el statement independiente.

\subsection{Varios candidatos mínimos}

El caso \texttt{many-minimums} consiste en un bloque \texttt{do} formado por dos cadenas independientes de tres statements cada una. El orden fuente comienza la primera cadena, desarrolla por completo la segunda y deja al final el último statement de la primera, de modo que distintos intercalados pueden exponer agrupaciones aplicativas equivalentes. Esta forma permite observar cómo la solución resuelve un empate entre varios reordenamientos válidos de costo mínimo. A continuación, en el Código~\ref{lst:resultados-many-minimums} y la Figura~\ref{fig:resultados-grafo-many-minimums} se muestran el programa probabilístico asociado a este caso y el grafo de precedencia RAW construido por la extensión, respectivamente.

\newpage

\begin{lstlisting}[style=haskell,caption={Programa probabilístico asociado al caso \texttt{many-minimums}.},label={lst:resultados-many-minimums}]
CD.do
  x1 <- Dist.uniform [1, 2]
  x2 <- Dist.uniform [x1 + 10, x1 + 20]
  x3 <- Dist.uniform [100, 200]
  x4 <- Dist.uniform [x3 + 1000, x3 + 2000]
  x5 <- Dist.uniform [x4 + 10000, x4 + 20000]
  x6 <- Dist.uniform [x2 + 100000, x2 + 200000]
  CD.return (x1, x2, x3, x4, x5, x6)
\end{lstlisting}

\begin{figure}[ht]
\centering
\begingroup
\setgraphnodesize{1.35cm}{0.9cm}{8pt}
\begin{PrecedenceGraph}[node distance=0.35cm and 0.28cm]
    \PrecedenceNode{s1}{x1}{$s_1$}
    \PrecedenceNode[right=of s1]{s2}{x2}{$s_2$}
    \PrecedenceNode[right=of s2]{s3}{x3}{$s_3$}
    \PrecedenceNode[right=of s3]{s4}{x4}{$s_4$}
    \PrecedenceNode[right=of s4]{s5}{x5}{$s_5$}
    \PrecedenceNode[right=of s5]{s6}{x6}{$s_6$}
    \PrecedenceEdge{s1}{s2}
    \PrecedenceEdge[bend left=38]{s2}{s6}
    \PrecedenceEdge{s3}{s4}
    \PrecedenceEdge{s4}{s5}
\end{PrecedenceGraph}
\endgroup
\caption[Grafo RAW del caso \texttt{many-minimums}.]{Grafo de precedencia RAW construido por la solución para el caso \texttt{many-minimums}.}
\label{fig:resultados-grafo-many-minimums}
\end{figure}

La Figura~\ref{fig:resultados-grafo-many-minimums} contiene las cadenas de aristas $s_1\to s_2\to s_6$ y $s_3\to s_4\to s_5$, sin dependencias entre ellas. Sus intercalados producen 20 órdenes válidos: 14 alcanzan el costo mínimo de 3 y los otros seis obtienen un costo de 4. La ejecución secuencial obtuvo un costo de 6 y \texttt{ApplicativeDo} redujo ese valor a 4 sobre el orden fuente. La solución obtuvo un costo de 3 al disponer cada cadena como una rama completa y seleccionó de manera estable el primero de los 14 candidatos mínimos. En conjunto, los cuatro perfiles confirman que la política funciona con mínimos únicos o empatados y no presupone que todos los casos deban mejorar el plan de \texttt{ApplicativeDo}.

% !TEX root = ../../main.tex

\section{Control negativo con IO}
\label{sec:resultados-io}

La evidencia positiva presentada en las secciones anteriores se obtuvo sobre programas evaluados bajo una hipótesis explícita de conmutatividad. Para contrastar ese escenario, esta sección retoma el control negativo con \texttt{IO} definido en el capítulo anterior. Su propósito es mostrar que un reordenamiento puede respetar todas las dependencias RAW y, aun así, modificar el comportamiento observable cuando la mónada utilizada no es conmutativa.

El programa sobre el que se aplicará el control negativo de la solución declara deliberadamente una instancia \texttt{CommutativeMonad IO} que no satisface la propiedad que esta clase representa. Dentro del bloque \texttt{CD.do}, el primer statement imprime \texttt{A} y produce el valor 1, mientras que el segundo imprime \texttt{B} y produce el valor 2. Ninguna acción lee el binder definido por la otra, pero ambas realizan efectos visibles sobre la salida estándar. A continuación, en el Código~\ref{lst:resultados-control-io} y la Figura~\ref{fig:resultados-grafo-control-io} se muestran el programa asociado a este control y el grafo de precedencia RAW construido por la extensión, respectivamente.

\begin{lstlisting}[style=haskell,caption={Programa asociado al control negativo con \texttt{IO}.},label={lst:resultados-control-io}]
instance CD.CommutativeMonad IO

program :: IO Int
program = CD.do
  x <- putStrLn "A" >> pure 1
  y <- putStrLn "B" >> pure 2
  CD.return (x + y)
\end{lstlisting}

\begin{figure}[ht]
\centering
\begin{PrecedenceGraph}
    \PrecedenceNode{s1}{x}{$s_1$}
    \PrecedenceNode[right=of s1]{s2}{y}{$s_2$}
\end{PrecedenceGraph}
\caption[Grafo RAW del control con \texttt{IO}.]{Grafo de precedencia RAW construido por la solución para el control negativo con \texttt{IO}.}
\label{fig:resultados-grafo-control-io}
\end{figure}

El grafo de la Figura~\ref{fig:resultados-grafo-control-io} no contiene aristas, ya que ninguno de los dos statements lee el binder producido por el otro. Por esta razón, el grafo admite los órdenes \texttt{[1,2]} y \texttt{[2,1]}, y ambos candidatos alcanzaron el costo mínimo. La ejecución secuencial obtuvo un costo de 2, mientras que \texttt{ApplicativeDo} y la solución obtuvieron un costo de 1. Como los dos candidatos empataron, la selección estable escogió el candidato de índice 0 y conservó el orden fuente \texttt{[1,2]}.

Desde la perspectiva del grafo y del modelo de costo, los dos órdenes son indistinguibles. Sin embargo, la ausencia de dependencias RAW no asegura que los efectos de las acciones puedan intercambiarse. Además del candidato seleccionado, la validación exhaustiva compiló y ejecutó el candidato forzado de índice 1, correspondiente al orden \texttt{[2,1]}. Ambos terminaron con código de salida cero y retornaron el entero 3, pero produjeron las secuencias observables resumidas en la Tabla~\ref{tab:resultados-control-io}.

\begin{table}[ht]
\centering
\TableSetup
\TableFrame{%
\begin{tabular}{lccc}
\rowcolor{tableHeaderBg}
\TableHead{Orden} & \TableHead{Primera línea} & \TableHead{Segunda línea} & \TableHead{Resultado} \\
\texttt{[1,2]} & \texttt{A} & \texttt{B} & \texttt{3} \\
\texttt{[2,1]} & \texttt{B} & \texttt{A} & \texttt{3} \\
\end{tabular}%
}
\caption{Salidas observables producidas por los dos órdenes del control negativo con \texttt{IO}.}
\label{tab:resultados-control-io}
\end{table}

Como muestra la Tabla~\ref{tab:resultados-control-io}, el orden fuente imprimió primero \texttt{A} y luego \texttt{B}, mientras que el candidato invertido produjo ambas líneas en el orden contrario. La comparación con el programa original detectó esta diferencia y terminó la validación con fracaso. Este resultado es el desenlace esperado del control y no representa un error de la implementación: confirma que un ordenamiento topológico preserva las dependencias locales de datos, pero no garantiza que los efectos de una mónada arbitraria puedan intercambiarse. La instancia \texttt{CommutativeMonad} expresa un contrato semántico que el compilador no demuestra, por lo que una declaración incorrecta puede habilitar transformaciones estructuralmente válidas que alteran el comportamiento observable.

% !TEX root = ../../main.tex

\section{Discusión}
\label{sec:resultados-discusion}

Respecto de \textbf{PV1}, los 93 programas sometidos al flujo de validación de la solución mostrado en la Figura~\ref{fig:validacion-flujo} conservaron el código de salida y la salida observable en \texttt{ApplicativeDo}, en la selección automática y en cada reordenamiento válido forzado. La coincidencia abarcó las familias estructurales de \texttt{Maybe} y \texttt{Dist.T Rational}, además de pesos no uniformes, soportes dependientes y condicionamiento probabilístico. Los casos analizados en la Sección~\ref{sec:resultados-casos-destacables} muestran que esta preservación se mantuvo ante grafos y construcciones sintácticas diferentes. La evidencia permite responder favorablemente la pregunta para los programas estudiados bajo la hipótesis explícita de conmutatividad, pero continúa siendo observacional y no constituye una demostración de equivalencia semántica para todo programa posible.

Respecto de \textbf{PV2}, en los mismos 93 programas la selección automática coincidió con el primer mínimo recalculado desde los costos de los candidatos compilados individualmente. La Tabla~\ref{tab:resultados-perfiles-costo} y sus cuatro casos muestran que la política se comportó correctamente cuando todos los candidatos empataron, cuando el orden original no fue mínimo, cuando existió un único mínimo y cuando varios reordenamientos compartieron el menor costo. Que un programa no reduzca el costo de \texttt{ApplicativeDo} no representa un fallo: PV2 exige seleccionar un mínimo entre los planes generados, no producir necesariamente una mejora estricta. Como el orden original forma parte del espacio explorado, el costo elegido tampoco puede superar el plan que \texttt{ApplicativeDo} obtiene sobre ese orden.

La cantidad de mejoras estrictas observadas no permite estimar la utilidad de la extensión en programas reales. Los corpus son sintéticos y la mayoría de sus casos fue diseñada para aislar preservación observable, dependencias, binders o empates, no para construir un orden fuente desfavorable. Las cinco variantes que sí mejoraron fueron replicadas con ambas mónadas y preparadas precisamente para admitir un plan menor. El resultado relevante no es su proporción dentro del corpus, sino que la extensión reconoció la alternativa mínima cuando existía y mantuvo un mínimo cuando no había una alternativa mejor.

Respecto de \textbf{PV3}, los seis controles distinguieron correctamente las rutas estándar y conmutativa. Con \texttt{ApplicativeDo} habilitado, los cuatro bloques \texttt{do} ordinarios mantuvieron en las trazas del Renamer \texttt{commutative-do = False}, aun en presencia de \texttt{QualifiedDo}, la importación de la interfaz y la instancia correspondiente. Los dos bloques \texttt{CD.do} produjeron \texttt{commutative-do = True} y generaron seis candidatos cada uno, uno original y cinco alternativos. Estos resultados confirman la activación local mediante el marcador dentro de las configuraciones evaluadas. La instancia \texttt{CommutativeMonad} sigue siendo necesaria para utilizar las operaciones calificadas, pero expresa una restricción comprobada durante el tipado y no constituye por sí sola el disparador del Renamer.

Los grafos presentados también muestran que el tamaño del espacio de búsqueda y la calidad del plan son propiedades relacionadas, pero diferentes. Una cadena completa produjo un único candidato; un grafo vacío de cuatro vértices produjo 24; y dos cadenas independientes de longitud cuatro produjeron 70. Sin embargo, el diamante mantuvo el mínimo que \texttt{ApplicativeDo} ya alcanzaba, mientras el orden desfavorable de las dos cadenas permitió reducir el costo de 6 a 4. Del mismo modo, las familias de \texttt{let}, binders, distribuciones dependientes, condicionamiento y \textit{shadowing} confirmaron que las aristas se forman a partir de las identidades y lecturas locales detectadas por el Renamer, no solo a partir de binders monádicos simples.

Finalmente, el control con \texttt{IO} de la Sección~\ref{sec:resultados-io} separa la validez estructural de la equivalencia semántica. El candidato invertido preservó el grafo RAW vacío y retornó el mismo entero, pero cambió el orden de los efectos visibles. Su divergencia no contradice los resultados positivos, sino que confirma el límite previsto por el diseño: las precedencias RAW protegen el flujo local de datos, mientras la posibilidad de intercambiar efectos depende de la conmutatividad declarada. Asimismo, las reducciones informadas corresponden al modelo estructural uniforme y no a mediciones de tiempo de ejecución. Bajo estas condiciones, los resultados responden favorablemente las tres preguntas de validación dentro del alcance delimitado en la Sección~\ref{sec:validacion-alcance}.

